\usepackage[spanish,es-nodecimaldot,es-tabla]{babel}
\usepackage[utf8]{inputenc}
\usepackage[T1]{fontenc}
\usepackage{lmodern}
\usepackage[a4paper,margin=2.5cm]{geometry}
\usepackage{amsmath,amssymb,amsthm,mathtools,bm}
\usepackage{graphicx}
\usepackage{booktabs,tabularx}
\usepackage{xcolor}
\usepackage[hidelinks]{hyperref}
\usepackage[style=apa,backend=biber]{biblatex}
\addbibresource{referencias.bib}

\graphicspath{{figuras/}}
% Mientras el notebook no haya generado una figura, se muestra un recuadro en su lugar.
\let\includegraphicsoriginal\includegraphics
\renewcommand{\includegraphics}[2][]{%
  \IfFileExists{figuras/#2.pdf}{\includegraphicsoriginal[#1]{#2}}{%
    \fbox{\parbox[c][4cm][c]{0.8\linewidth}{\centering Figura pendiente:\\ \texttt{\detokenize{#2}}}}}}

\theoremstyle{plain}
\newtheorem{teorema}{Teorema}[chapter]
\newtheorem{proposicion}[teorema]{Proposición}
\newtheorem{lema}[teorema]{Lema}
\theoremstyle{definition}
\newtheorem{definicion}[teorema]{Definición}
\newtheorem{hipotesis}{Hipótesis}
\renewcommand{\thehipotesis}{H\arabic{hipotesis}}
\renewcommand{\qedsymbol}{$\blacksquare$}
\renewcommand{\proofname}{Demostración}

\newcommand{\R}{\mathbb{R}}
\newcommand{\E}{\mathbb{E}}
\newcommand{\simplex}{\Delta^{N-1}}
\newcommand{\TPR}{\mathrm{TPR}}
\newcommand{\FPR}{\mathrm{FPR}}
\newcommand{\pendiente}[1]{\textcolor{orange}{[\textbf{Pendiente:} #1]}}
